\appendix

\section{Derivation of Gyroscope Integration Error Bounds}
\label{app:error_bounds}

This appendix provides a rigorous derivation of the three error terms that
arise when numerically integrating a gyroscope signal: bias drift, angular
random walk, and discretisation error. These results underpin the theoretical
envelopes shown in Figs.~\ref{fig:gyro_drift} and~\ref{fig:gyro_sampling}.

% ===================================================================
\subsection{Setup and Notation}
\label{app:setup}

Consider a single-axis gyroscope sampled at frequency $f_s = 1/\Delta t$.
At each time step $n$, the measured angular velocity is:

\begin{equation}
  \omega[n] = \omega_\text{true}[n] + b_\omega + n_\omega[n]
  \label{eq:app:gyro_model}
\end{equation}

\noindent where:
\begin{itemize}
  \item $\omega_\text{true}[n]$ is the true angular velocity,
  \item $b_\omega$ is a constant bias (deg/s),
  \item $n_\omega[n]$ are independent and identically distributed (i.i.d.)
        random variables with $E[n_\omega[n]] = 0$ and
        $\text{Var}[n_\omega[n]] = \sigma_\omega^2$.
\end{itemize}

The angle is estimated by rectangular (Euler) integration:

\begin{equation}
  \hat{\theta}[N] = \sum_{n=1}^{N} \omega[n]\,\Delta t
  \label{eq:app:integration}
\end{equation}

\noindent while the true angle is:

\begin{equation}
  \theta_\text{true}[N] = \sum_{n=1}^{N} \omega_\text{true}[n]\,\Delta t
  \label{eq:app:true_angle}
\end{equation}

The total error is the difference:

\begin{equation}
  \epsilon[N] = \hat{\theta}[N] - \theta_\text{true}[N]
              = \sum_{n=1}^{N} \bigl(b_\omega + n_\omega[n]\bigr)\,\Delta t
  \label{eq:app:total_error}
\end{equation}

\noindent We can separate this into two independent terms:

\begin{equation}
  \epsilon[N]
  = \underbrace{\sum_{n=1}^{N} b_\omega\,\Delta t}_{\epsilon_\text{bias}[N]}
  + \underbrace{\sum_{n=1}^{N} n_\omega[n]\,\Delta t}_{\epsilon_\text{rw}[N]}
  \label{eq:app:error_split}
\end{equation}

% ===================================================================
\subsection{Bias Drift: $\epsilon_\text{bias}$}
\label{app:bias}

Since $b_\omega$ is constant, the sum is trivial:

\begin{equation}
  \epsilon_\text{bias}[N]
  = \sum_{n=1}^{N} b_\omega\,\Delta t
  = N \cdot b_\omega \cdot \Delta t
  \label{eq:app:bias_sum}
\end{equation}

\noindent Substituting $T = N\,\Delta t$ (total elapsed time):

\begin{equation}
  \boxed{\;\epsilon_\text{bias} = b_\omega \cdot T\;}
  \label{eq:app:bias_final}
\end{equation}

\paragraph{Properties.}
\begin{itemize}
  \item \textbf{Deterministic}: the bias drift is the same in every
        realisation of the experiment.
  \item \textbf{Linear in time}: doubling the measurement duration doubles
        the error.
  \item \textbf{Independent of $\Delta t$}: for a fixed total time $T$,
        changing the sampling frequency does not affect the bias drift.
        This is because $N \cdot \Delta t = T$ regardless of how we
        partition the interval.
\end{itemize}

% ===================================================================
\subsection{Angular Random Walk: $\epsilon_\text{rw}$}
\label{app:rw}

This is the stochastic term and requires a careful probabilistic analysis.

\subsubsection{Step 1: Write the sum explicitly}

\begin{equation}
  \epsilon_\text{rw}[N]
  = \sum_{n=1}^{N} n_\omega[n]\,\Delta t
  = \Delta t \sum_{n=1}^{N} n_\omega[n]
  \label{eq:app:rw_sum}
\end{equation}

\subsubsection{Step 2: Compute the expected value (mean)}

Since $E[n_\omega[n]] = 0$ for all $n$, and expectation is linear:

\begin{align}
  E\bigl[\epsilon_\text{rw}[N]\bigr]
  &= \Delta t \sum_{n=1}^{N} E[n_\omega[n]] \notag \\
  &= \Delta t \sum_{n=1}^{N} 0 \notag \\
  &= 0
  \label{eq:app:rw_mean}
\end{align}

\noindent The random walk has zero mean---it does not favour any direction on
average.

\subsubsection{Step 3: Compute the variance}

The variance of a sum of \emph{independent} random variables equals the sum of
their variances. This is the key property we exploit. Since $n_\omega[n]$ are
i.i.d.:

\begin{align}
  \text{Var}\bigl[\epsilon_\text{rw}[N]\bigr]
  &= \text{Var}\!\left[\Delta t \sum_{n=1}^{N} n_\omega[n]\right] \notag \\[6pt]
  &= \Delta t^2 \;\text{Var}\!\left[\sum_{n=1}^{N} n_\omega[n]\right]
  \label{eq:app:rw_var_step1}
\end{align}

\noindent where we used the property $\text{Var}[cX] = c^2\,\text{Var}[X]$ to
pull out the constant $\Delta t^2$. Now, for independent variables:

\begin{align}
  \text{Var}\!\left[\sum_{n=1}^{N} n_\omega[n]\right]
  &= \sum_{n=1}^{N} \text{Var}[n_\omega[n]] \notag \\[4pt]
  &= \sum_{n=1}^{N} \sigma_\omega^2 \notag \\[4pt]
  &= N\,\sigma_\omega^2
  \label{eq:app:rw_var_step2}
\end{align}

\noindent Substituting into Eq.~\eqref{eq:app:rw_var_step1}:

\begin{equation}
  \text{Var}\bigl[\epsilon_\text{rw}[N]\bigr]
  = \Delta t^2 \cdot N \cdot \sigma_\omega^2
  = N\,\Delta t^2\,\sigma_\omega^2
  \label{eq:app:rw_variance}
\end{equation}

\subsubsection{Step 4: Compute the standard deviation}

The standard deviation (the ``1-sigma envelope'') is the square root of the
variance:

\begin{equation}
  \sigma_{\epsilon_\text{rw}}
  = \sqrt{N\,\Delta t^2\,\sigma_\omega^2}
  = \sigma_\omega\,\Delta t\,\sqrt{N}
  \label{eq:app:rw_std_N}
\end{equation}

\subsubsection{Step 5: Express in terms of physical quantities}

Substituting $N = T / \Delta t$:

\begin{align}
  \sigma_{\epsilon_\text{rw}}
  &= \sigma_\omega\,\Delta t\,\sqrt{\frac{T}{\Delta t}} \notag \\[6pt]
  &= \sigma_\omega\,\Delta t \cdot \frac{\sqrt{T}}{\sqrt{\Delta t}} \notag \\[6pt]
  &= \sigma_\omega\,\sqrt{\Delta t}\,\sqrt{T} \notag \\[6pt]
  &= \sigma_\omega\,\sqrt{\Delta t \cdot T}
  \label{eq:app:rw_std_derivation}
\end{align}

\begin{equation}
  \boxed{\;\sigma_{\epsilon_\text{rw}}
  = \sigma_\omega\,\sqrt{\Delta t \cdot T}\;}
  \label{eq:app:rw_final}
\end{equation}

\paragraph{Interpretation.}
\begin{itemize}
  \item The standard deviation grows as $\sqrt{T}$---much slower than the
        linear bias drift. This is the hallmark of a \emph{random walk}.
  \item Unlike the bias drift, the random walk \emph{does} depend on
        $\Delta t$: a higher sampling frequency (smaller $\Delta t$) yields a
        \emph{smaller} random walk, because each noise sample is multiplied by
        a smaller time step before being accumulated.
  \item The $\pm 1\sigma$ envelope contains approximately 68\% of the
        realisations (assuming Gaussian noise). The $\pm 2\sigma$ envelope
        contains $\sim$95\%.
\end{itemize}

\subsubsection{Relation to the Angular Random Walk coefficient}

In sensor datasheets, manufacturers report the \emph{Angular Random Walk}
(ARW) coefficient, defined as:

\begin{equation}
  \sigma_\text{ARW} = \sigma_\omega\,\sqrt{\Delta t}
  = \frac{\sigma_\omega}{\sqrt{f_s}}
  \qquad \text{(units: deg/}\sqrt{\text{s}}\text{)}
  \label{eq:app:arw_def}
\end{equation}

\noindent so that Eq.~\eqref{eq:app:rw_final} simplifies to:

\begin{equation}
  \sigma_{\epsilon_\text{rw}} = \sigma_\text{ARW}\,\sqrt{T}
  \label{eq:app:arw_simple}
\end{equation}

\noindent This compact form is commonly used in inertial navigation literature.

% ===================================================================
\subsection{Discretisation Error}
\label{app:disc}

Even if the gyroscope were \emph{perfect} ($b_\omega = 0$,
$\sigma_\omega = 0$), a third error arises from approximating a continuous
integral with a discrete sum. We derive its order using Taylor expansion.

\subsubsection{Step 1: Local truncation error}

The true integral over one time step is:

\begin{equation}
  \int_{t_n}^{t_n + \Delta t} \omega(\tau)\,d\tau
  \label{eq:app:true_step}
\end{equation}

\noindent The rectangular (Euler) rule approximates this as
$\omega(t_n)\,\Delta t$. Expanding $\omega(\tau)$ in a Taylor series around
$t_n$:

\begin{equation}
  \omega(\tau) = \omega(t_n) + \dot{\omega}(t_n)(\tau - t_n)
              + \mathcal{O}\bigl((\tau - t_n)^2\bigr)
  \label{eq:app:taylor}
\end{equation}

\noindent Integrating from $t_n$ to $t_n + \Delta t$:

\begin{align}
  \int_{t_n}^{t_n + \Delta t} \omega(\tau)\,d\tau
  &= \omega(t_n)\,\Delta t
   + \dot{\omega}(t_n)\,\frac{\Delta t^2}{2}
   + \mathcal{O}(\Delta t^3) \notag \\[6pt]
  &= \underbrace{\omega(t_n)\,\Delta t}_{\text{Euler approx.}}
   + \underbrace{\dot{\omega}(t_n)\,\frac{\Delta t^2}{2}
   + \mathcal{O}(\Delta t^3)}_{\text{local truncation error}}
  \label{eq:app:local_error}
\end{align}

\noindent Therefore, the local truncation error per step is:

\begin{equation}
  e_\text{local} = \dot{\omega}(t_n)\,\frac{\Delta t^2}{2}
                 = \mathcal{O}(\Delta t^2)
  \label{eq:app:local_order}
\end{equation}

\subsubsection{Step 2: Global accumulated error}

Over $N = T/\Delta t$ steps, the errors accumulate:

\begin{align}
  \epsilon_\text{disc}
  &= \sum_{n=1}^{N} e_\text{local}[n] \notag \\[4pt]
  &= \sum_{n=1}^{N} \mathcal{O}(\Delta t^2) \notag \\[4pt]
  &= N \cdot \mathcal{O}(\Delta t^2) \notag \\[4pt]
  &= \frac{T}{\Delta t} \cdot \mathcal{O}(\Delta t^2)
  \label{eq:app:global_derivation}
\end{align}

\begin{equation}
  \boxed{\;\epsilon_\text{disc} = \mathcal{O}(\Delta t)\;}
  \label{eq:app:disc_final}
\end{equation}

\paragraph{Interpretation.}
\begin{itemize}
  \item The global discretisation error is \textbf{first-order} in $\Delta t$.
        Halving the time step (doubling the sampling frequency) halves this
        error.
  \item For a sensor sampling at $f_s = 100$~Hz ($\Delta t = 0.01$~s), this
        error is typically very small compared to the bias drift, which is why
        increasing the sampling rate has a marginal effect in practice.
\end{itemize}

% ===================================================================
\subsection{Complete Error Budget}
\label{app:budget}

Combining all three terms, the total angle error after time $T$ is:

\begin{equation}
  \epsilon_\text{total}(T)
  = \underbrace{b_\omega \cdot T\vphantom{\sqrt{}}}_{\text{bias drift}}
  + \underbrace{\sigma_\omega\sqrt{\Delta t \cdot T}}_{\text{random walk (stochastic)}}
  + \underbrace{\mathcal{O}(\Delta t)}_{\text{discretisation}}
  \label{eq:app:total_budget}
\end{equation}

\noindent Table~\ref{tab:app:error_comparison} compares the scaling of each
term with respect to both elapsed time~$T$ and sampling interval~$\Delta t$.

\begin{table}[h]
  \centering
  \caption{Scaling of the three gyroscope integration error terms.}
  \label{tab:app:error_comparison}
  \begin{tabular}{@{}lccc@{}}
    \toprule
    \textbf{Error term}
      & \textbf{Grows with $T$}
      & \textbf{Depends on $\Delta t$?}
      & \textbf{Dominant?} \\
    \midrule
    Bias drift
      & $\propto T$ (linear)
      & No
      & Yes (long term) \\
    Random walk
      & $\propto \sqrt{T}$
      & Yes ($\propto \sqrt{\Delta t}$)
      & Medium \\
    Discretisation
      & Fixed for fixed $T$
      & Yes ($\propto \Delta t$)
      & Negligible \\
    \bottomrule
  \end{tabular}
\end{table}

\paragraph{Numerical example.}
For a typical MEMS gyroscope with $b_\omega = 0.05$~deg/s,
$\sigma_\omega = 0.1$~deg/s, at $f_s = 100$~Hz ($\Delta t = 0.01$~s):

\begin{center}
\begin{tabular}{@{}lcccc@{}}
  \toprule
  & $T = 1$~s & $T = 10$~s & $T = 60$~s & $T = 300$~s \\
  \midrule
  $\epsilon_\text{bias}$
    & $0.05°$ & $0.50°$ & $3.00°$ & $15.00°$ \\
  $\sigma_{\epsilon_\text{rw}}$
    & $0.01°$ & $0.03°$ & $0.08°$ & $0.17°$ \\
  Ratio (bias/rw)
    & $5\times$ & $16\times$ & $39\times$ & $87\times$ \\
  \bottomrule
\end{tabular}
\end{center}

\noindent The bias drift dominates increasingly with time: after 5~minutes the
bias error is 87 times larger than the random walk. This justifies the use of
external corrections (complementary filter, ZUPT) to bound the drift.
